\documentclass[12pt]{article}
\usepackage[a4paper, margin=2cm]{geometry}  % マージンを狭くする
\usepackage{setspace}  % 行間調整
\setstretch{1.5}  % 行間を広げる

\begin{document}

\begin{center}
\textbf{修士論文予備審査報告書に対する回答書}  % 太字にしたい場合
\end{center}
(1)画像処理における点の拡大率がクラスター数などに影響を与えるが、点の拡大率を自動的に選択することができるか検討すること。\\\\
データを画像化する目的は、データの輪郭を基に細線化を適用し、大まかな木構造を計算することである。この木構造の葉と分岐の数をクラスタ数として用いる。
解像度が低すぎるとデータの輪郭を正確に捉えられず、高すぎると画像内の連結オブジェクトの数が増加し、結果としてクラスタ数が過大になる可能性がある。
具体的な手法として、初期解像度はデータの分散の三乗根を基準とし、初期解像度の10倍までの範囲で9段階の解像度を設定して画像化を行う。各解像度において、画像の連結オブジェクト数とオブジェクトの面積を算出する。その後、データの凸包面積以下であり、かつ連結オブジェクト数が10〜50の範囲にある解像度の中から、オブジェクトの面積が最大となる解像度を選択する（節2.2）。
この方法により、適切な解像度を自動的に決定し、解像度が粗すぎず細かすぎないバランスの取れた画像化が可能となる。
\\
\\
(2)研究の目的や評価結果を明確にすること。特に、評価法は独自のものではなく他の研究で使用されるものを適用できないかを検討し、複数のデータを用いて既存手法との十分な比較を行うこと。\\\\
研究の目的は、細胞系統擬似時間解析において、画像処理によるパラメータの自動決定と細胞状態が3つ以上に分岐しないような制約を追加することによって予測精度を向上させることである(章1)。
データはシミュレーションデータではなく、公開データを4つ用意し比較を行った(章4)。
評価方法は他の研究で使用されるものを適用できないかを検討したが、制約による予測精度向上を検証したいため公開データを使用して制約ありなしの解析結果を比較するにとどめた。
制約を追加することで細胞状態の分岐が3つ以上に遷移しなくなることを確認した(節3.1)。
既存手法であるMonocleでの解析結果と制約ありなしの結果を比較し、制約ありの結果がMonocleの解析結果と近くなることを確認した(節3.3)。
DNAバーコードは、特定の DNA 領域にランダムな挿入するバーコードのことであり細胞が分裂するたびに引き継ぐことにより細胞の系統的な起源を識別することができるものである。
予測精度向上については、DNAバーコード付きのデータを使用して細胞の兄弟関係の予測精度を制約ありなしで比較した(節3.4)。

\end{document}
